\documentclass[10pt,a4paper]{article}
\usepackage[utf8]{inputenc}
\usepackage[spanish,es-nodecimaldot]{babel}
\usepackage{amsmath,amssymb,geometry,hyperref,fancyhdr,microtype,booktabs,float}
\geometry{margin=2cm}\pagestyle{fancy}\fancyhf{}
\fancyhead[L]{\small Universidad Externado de Colombia}\fancyhead[R]{\small M\'etodos Numéricos -- Proyecto Integrador (Fase 2)}\fancyfoot[C]{\thepage}
\hypersetup{colorlinks=true,linkcolor=black,citecolor=black,urlcolor=blue}
\begin{document}
\begin{titlepage}\centering\vspace*{2cm}{\LARGE\bfseries Proyecto Integrador: Segunda Entrega}\\[.5cm]{\large Implementación reproducible para ra\'ices m\'ultiples}\\[.6cm]{\Large M\'etodos Numéricos}\\[2cm]Integrantes: Ailyn Sof\'ia Gómez Rodríguez, Juli\'an Camilo Jiménez, Juli\'an Steven Duarte\\[1cm]Docente: Edwin Fernando Mu\~noz Chipatecua\vfill Pregrado en Ciencia de Datos\\Departamento de Matemáticas\\Universidad Externado de Colombia\\Bogotá D.C., Colombia -- Noviembre de 2026\end{titlepage}
\section{El problema que estudiamos}
Resolver una ecuaci\'on no lineal $g(x)=0$ significa encontrar un valor de $x$ para el cual la funci\'on vale cero. A veces, la ra\'iz aparece repetida: se dice que tiene multiplicidad $m$. Cerca de esa ra\'iz la funci\'on se aplana ($g'(x)\to 0$), provocando que el m\'etodo de Newton cl\'asico pierda su convergencia cuadr\'atica y decaiga a orden lineal ($p=1$), mientras que los errores de redondeo amplifican la cancelaci\'on num\'erica \cite{turner}. Si se conoce $m$, Newton modificado recupera el orden cuadr\'atico ($p=2$). El art\'iculo de Jerezano, Chicharro y Garrido-S\'aez (2026) presenta m\'etodos que aprovechan el conocimiento de $m$ y evitan calcular derivadas anal\'iticas \cite{jerezano}. Nuestro objetivo es implementar sus dos variantes principales, validarlas experimentalmente frente a l\'ineas base del curso y analizar sus implicaciones para problemas de optimizaci\'on y ciencia de datos.
\section{M\'etodos num\'ericos}
Una diferencia dividida calcula la pendiente entre dos evaluaciones: $g[w_k,x_k]=(g(w_k)-g(x_k))/(w_k-x_k)$. La familia del art\'iculo primero aproxima esa pendiente y calcula un punto intermedio:
\begin{align*}w_k&=x_k+\beta g(x_k)^2,& y_k&=x_k-m\frac{g(x_k)}{g[w_k,x_k]},\\t_k&=\left(\frac{g(y_k)}{g(x_k)}\right)^{1/m},&x_{k+1}&=x_k-mH(t_k)\frac{g(x_k)}{g[w_k,x_k]} .\end{align*}
La potencia fraccionaria se calcula empleando la rama principal del logaritmo complejo, es decir, $t_k = \exp\left(\frac{1}{m}\mathrm{Ln}\left(\frac{g(y_k)}{g(x_k)}\right)\right)$, donde $\mathrm{Ln}(z) = \ln|z| + i\operatorname{Arg}(z)$ con $\operatorname{Arg}(z)\in(-\pi,\pi]$. Esta formulaci\'on metodol\'ogica es fundamental y te\'oricamente rigurosa, ya que previene discontinuidades o indefiniciones num\'ericas durante el c\'alculo de potencias fraccionarias cuando el cociente $g(y_k)/g(x_k)$ toma argumentos negativos o valores complejos a lo largo de la trayectoria iterativa, garantizando que el m\'etodo mantenga una evaluaci\'on univaluada, continua y anal\'itica en todo el plano complejo.
Se estudian $H_1(t)=1+t+2t^2$ (M1) y $H_2(t)=(1-t)/(1-2t)$ (M2), con $\beta=0.01$. Bajo las condiciones del art\'iculo, ambos tienen orden te\'orico cuatro y requieren tres evaluaciones funcionales. Como l\'ineas base implementamos Newton modificado, $x_{k+1}=x_k-mg(x_k)/g'(x_k)$, y Steffensen modificado, que aproxima la derivada con una secante. Ambos tienen orden te\'orico dos para ra\'ices m\'ultiples. También implementamos el m\'etodo KK, competidor libre de derivadas de orden cuatro que aparece en las tablas del art\'iculo. En KK se usa $w_k=x_k+\beta g(x_k)$, el mismo predictor $y_k$, y el corrector $x_{k+1}=y_k-((m+2)t_k/(1-2t_k))g(x_k)/(g[w_k,x_k]+2g[w_k,y_k])$.
El orden te\'orico describe cómo se reduce el error cuando ya estamos cerca de la ra\'iz. En la práctica tambi\'en importa si la semilla está en una zona desde la cual el m\'etodo converge, si las divisiones están bien definidas y si el redondeo impide seguir avanzando.
\section{Funciones y dise\~no experimental}
Usamos las funciones deterministas descritas en el art\'iculo: el polinomio cúbico $g_1(x)=x^3-5.22x^2+9.0825x-5.2675$, con ra\'iz doble objetivo $1.75$ y ra\'iz simple cercana $1.72$; la funci\'on $g_2(x)=(e^{-x}-1+x/5)^3$, correspondiente a la ecuaci\'on de Planck al cubo, con ra\'iz triple aproximada $4.96511423$; y el polinomio $g_3$ de grado nueve definido en el notebook, con una ra\'iz de multiplicidad cuatro en $3$. La expresión de Planck es la ecuaci\'on base elevada al cubo; el término lineal es x sobre 5, no x al cubo sobre 5.
En g2 se estudia la raíz positiva distinta de cero. La función también se anula en cero, pero esa raíz no es la que se analiza con las semillas escogidas.
Para cada problema se prueban dos semillas iniciales fijadas en el cuaderno. Esto permite comparar más de una trayectoria y observar fallos de convergencia. Las ra\'ices conocidas se usan solo para calcular el error absoluto; los iteradores no reciben la solución. Las funciones son datos matemáticos reproducibles, por lo que no hay muestreo aleatorio ni semilla pseudoaleatoria que controlar. Además de la corrida base con $\beta=0.01$, el notebook repite M1 y M2 con $\beta$ en $\{0.001,0.01,1.0\}$ y muestra una prueba complementaria en doble precisi\'on para observar estancamiento.
\section{Implementación reproducible}
El notebook se organiza en configuración, definición de funciones, solucionadores, ejecución de corridas, tabla, gr\'afico y lectura de resultados. Requiere Python 3 y los paquetes mpmath, pandas y matplotlib. Se fijan precisi\'on decimal arbitraria, tolerancia, n\'umero m\'aximo de iteraciones, beta y semillas. Las constantes decimales se crean desde texto para evitar redondeo binario innecesario. La precisi\'on alta ayuda a reducir absorción por redondeo, pero no asegura convergencia.
Para la potencia fraccionaria se usa la rama principal compleja, como en la primera entrega; las trayectorias pueden ser complejas aunque las semillas y raíces objetivo sean reales. La corrida termina cuando paso y residuo son menores que la tolerancia. Se registran estancamiento, singularidades y límite de iteraciones.
\section{Criterios de comparaci\'on y validaci\'on experimental del orden (ACOC)}
Se guardan el n\'umero de iteraciones, error absoluto $|x_k-\alpha|$, residuo, ACOC y tiempo transcurrido. Siguiendo la ruta pedag\'ogica del curso (Teor\'ia $\to$ Implementaci\'on $\to$ An\'alisis $\to$ Aplicaci\'on) expuesta por Turner et al. \cite{turner}, el ACOC estima computacionalmente el orden asint\'otico usando tres diferencias sucesivas:
\[\mathrm{ACOC}\approx\frac{\log(|x_{k+1}-x_k|/|x_k-x_{k-1}|)}{\log(|x_k-x_{k-1}|/|x_{k-1}-x_{k-2}|)}.\]
Para validar la implementaci\'on del ACOC previo al an\'alisis de multiplicidad, el cuaderno eval\'ua los m\'etodos can\'onicos de la Semana 9 sobre $f(x)=x^3-x-2=0$ (ra\'iz simple $\alpha\approx 1.5213797$): Newton converge en 7 pasos hacia su orden cuadr\'atico te\'orico con $\mathrm{ACOC}\to 2.0000$, mientras la secante cl\'asica converge en 9 pasos hacia el orden superlineal te\'orico $p=\frac{1+\sqrt{5}}{2}\approx 1.6180$ con $\mathrm{ACOC}\to 1.6162$.

En ra\'ices m\'ultiples, el ACOC solo es significativo cuando el m\'etodo est\'a en su r\'egimen asint\'otico terminal y los pasos a\'un superan la precisi\'on finita. Por ello, el algoritmo restringe el ACOC a ejecuciones que cumplan la condici\'on estricta de convergencia y cuenten con al menos 4 iterados v\'alidos, evitando reportar valores espurios ante paradas prematuras. El tiempo depende del equipo; adem\'as, M1 y M2 eval\'uan tres puntos de $g$ por paso, de modo que menos iteraciones no implican autom\'aticamente menor costo total.
\section{Precisión, estabilidad, parada y límites num\'ericos}
En la familia sin derivadas, el nuevo punto $w_k$ se obtiene sumando a $x_k$ una cantidad proporcional a $g(x_k)^2$ ($o$ a $g(x_k)$ en Steffensen y KK). Al aproximarse a la ra\'iz, esa perturbaci\'on decae con extrema rapidez. En aritm\'etica finita, si $|w_k-x_k|$ es inferior a la unidad en el \'ultimo lugar (ULP), ocurre absorci\'on num\'erica ($w_k == x_k$), imposibilitando la construcci\'on del cociente secante $g[w_k,x_k]$. Nuestro c\'odigo detecta expl\'icitamente esta p\'erdida de perturbaci\'on y clasifica la salida como \texttt{'estancamiento w=x'} sin recurrir al artificio de cambiar a una derivada.

Asimismo, cuando la derivada anal\'itica se aproxima a cero ($g'(x_n)\approx 0$) o la pendiente secante se anula ($g[w_k,x_k]\approx 0$) por oscilaciones debidas a una semilla lejana o a un punto de inflexi\'on local, la f\'ormula iterativa genera un paso hiperb\'olico divergente ($x_{k+1}\to\infty$). Para asegurar estabilidad computacional, el solucionador captura excepciones \texttt{ZeroDivisionError('derivada nula')} y \texttt{ZeroDivisionError('pendiente secante nula')}, as\'i como cancelaciones catastr\'oficas en los denominadores de KK y polos de $H_2(t)$ ($1-2t=0$), deteniendo la corrida controladamente e informando el diagn\'ostico en la tabla de estados.

Se eligieron 150 d\'igitos de precisi\'on (mpmath) y tolerancia $10^{-25}$ como configuraci\'on de trabajo para analizar la din\'amica asint\'otica con amplio margen num\'erico frente a la doble precisi\'on binaria (\texttt{float64} IEEE-754). Para ra\'ices m\'ultiples, el residuo decae aproximadamente como $|x-\alpha|^m$; por ello, un residuo de $10^{-100}$ en $m=4$ solo garantiza un error del orden de $10^{-25}$, justificando la exigencia simult\'anea de tolerancia en residuo y en paso ($|x_{k+1}-x_k|<\mathrm{TOL}$ y $|g(x_{k+1})|<\mathrm{TOL}$).
\section{Resultados de la corrida de referencia}
El notebook se ejecutó con 150 d\'igitos decimales, tolerancia $10^{-25}$, m\'aximo de 100 iteraciones y $\beta=0.01$. El Cuadro~\ref{tab:resultados} resume casos seleccionados. Los estados de estancamiento se mantienen aunque el residuo sea peque\~no: el criterio tambi\'en exige que el paso cumpla la tolerancia.
\begin{table}[H]
\centering
\caption{Comparación de corridas representativas con precisi\'on de 150 d\'igitos y tolerancia $10^{-25}$.}
\label{tab:resultados}
\small
\begin{tabular}{llrrrrrl}
\toprule
Problema/semilla & Método & Iter. & $g$ eval. & $g'$ eval. & Error abs. & Residuo & Estado \\
\midrule
$g_1$, $x_0=1.9$ & M1 & 5 & 16 & 0 & $1.32\times10^{-35}$ & 5.26e-72 & Convergió \\
$g_1$, $x_0=1.9$ & M2 & 4 & 13 & 0 & $2.03\times10^{-50}$ & 1.24e-101 & Estancamiento \\
$g_1$, $x_0=1.9$ & Newton mod. & 8 & 9 & 8 & $1.69\times10^{-57}$ & 8.58e-116 & Convergió \\
$g_1$, $x_0=1.9$ & Steffensen & 8 & 17 & 0 & $1.77\times10^{-57}$ & 9.43e-116 & Convergió \\
$g_1$, $x_0=1.9$ & KK & 4 & 14 & 0 & $1.01\times10^{-53}$ & 3.06e-108 & Pendiente nula \\
$g_3$, $x_0=2.8$ & M1 & 3 & 10 & 0 & $3.71\times10^{-18}$ & 1.52e-68 & Convergió \\
$g_3$, $x_0=2.8$ & M2 & 3 & 10 & 0 & $5.42\times10^{-19}$ & 6.89e-72 & Convergió \\
$g_2$, $x_0=5.3$ & M1 & 2 & 7 & 0 & $2.07\times10^{-28}$ & $6.41\times10^{-86}$ & Estancamiento \\
$g_2$, $x_0=5.3$ & KK & 3 & 10 & 0 & $9.91\times10^{-117}$ & $7.01\times10^{-351}$ & Convergió \\
\bottomrule
\end{tabular}
\end{table}
Para $g_1$ con $x_0=1.9$, el ACOC observado es 8.81 para M1 (único de la familia que satisface ambos criterios de parada) y cercano a 2 para las dos l\'ineas base (Newton modificado y Steffensen con 2.00). En M2 y KK el ACOC no se calcula al terminar en estancamiento por absorción secante ($w=x$) o división por pendiente secante nula, respectivamente. El valor de M1 no coincide con el orden te\'orico (orden 4) y obliga a revisar la historia iterativa, no a afirmar prematuramente que se verificó dicho orden. El cálculo de ACOC en las últimas iteraciones es altamente sensible e inestable cuando los pasos sucesivos se aproximan a la escala de la tolerancia o a la saturación numérica de la máquina; por ello, en el código y en el análisis se restringe el ACOC a iteraciones exitosamente convergentes con al menos cuatro iterados, presentándolo estrictamente como una herramienta diagnóstica local. Los tiempos promedio reportados corresponden al promedio de 50 ejecuciones y dependen del hardware y entorno del sistema.
\section{Lectura comparativa de los resultados}
El Cuadro~\ref{tab:resultados} muestra que, para $g_1$ con $x_0=1.9$, M1 y M2 alcanzan aproximaciones con errores absolutos menores que $10^{-35}$. M1 satisface los dos criterios de parada en cinco iteraciones. M2 llega a un error aún menor, pero termina por estancamiento: el paso final no basta para cumplir la tolerancia aunque el residuo sea diminuto. Esta diferencia ilustra por qué se registran el estado, el residuo y el error, en vez de ordenar los m\'etodos solo por una cifra.
En el mismo caso, las dos l\'ineas base llegan a errores del orden de $10^{-57}$ en ocho iteraciones. Sus evaluaciones de $g'$ quedan registradas por separado; M1 y M2 no usan derivadas. El conteo de evaluaciones en el CSV incluye las necesarias para verificar el residuo y evita comparar solamente el n\'umero de iteraciones. El tiempo informado es el promedio de 50 ejecuciones de cada combinación y sigue siendo sensible al sistema, por lo que no se presenta como una ventaja general. KK tambi\'en se incluye en el CSV y en el notebook; en la corrida de referencia $g_1$, $x_0=1.9$, se detiene por pendiente secante nula antes de cumplir la tolerancia.
Las corridas restantes confirman que el resultado depende de la semilla: para g1 desde 3.0 convergen las dos líneas base y KK; para g2 KK converge desde 5.3, pero no desde 10.0. En g3 convergen todos los métodos; M1 y M2 requieren tres o cuatro iteraciones y KK cuatro o seis. Con multiplicidad cuatro, el residuo puede ser mucho menor que el error, por lo que ambas medidas deben leerse juntas.
En la sensibilidad de $g_1$ con semilla $1.9$, M1 converge en cinco pasos para los tres valores de $\beta$, mientras M2 se estanca en los tres; este resultado local no permite afirmar que el par\'ametro no afecte otras semillas o funciones. En la prueba float64, M1 se absorbe en $w=x$ tras dos pasos y M2 encuentra una pendiente secante nula, con residuos finales del orden de $10^{-9}$ y $10^{-6}$, respectivamente.
El barrido también muestra que el efecto de beta no es monotónico. Para g1 desde 3.0, M1 falla con beta=0.001 y 0.01, pero converge con beta=1.0; M2 se estanca con los dos primeros valores y alcanza el límite de 100 iteraciones con beta=1.0. En g3, beta=1.0 tampoco produce convergencia con las semillas probadas. Estos casos muestran que el parámetro puede cambiar cada trayectoria, pero no permiten inferir una región general de convergencia.
El número de iteraciones tampoco basta para comparar el costo. En g1 con semilla 1.9, M1 usa 16 evaluaciones de g en cinco pasos; Newton usa 9 evaluaciones de g y 8 de su derivada en ocho pasos, y Steffensen usa 17 evaluaciones de g. El tiempo promedio también depende del equipo. Por eso presentamos estas cantidades por separado y no concluimos que M1 sea más eficiente solo porque termina en menos iteraciones.

En g2 con semilla 5.3, M1 y M2 se detienen cuando w coincide con x, aunque sus residuos son 6.41e-86 y 6.60e-87. KK, en cambio, converge en tres pasos con residuo 7.01e-351. Este caso muestra por qué distinguimos entre una aproximación muy cercana y una corrida que satisface el criterio de parada completo.
El ACOC de $g_1$ con semilla $1.9$ es 8.81 para M1 y exactamente 2.00 para las l\'ineas base (descartando estimaciones espurias en M2 y KK al no converger bajo la tolerancia dual). El valor observado en M1 no confirma el orden teórico cuatro: con pocas iteraciones (5 pasos), una medida local refleja la aceleración final del transitorio o efectos de escala numérica. Para $g_3$ también aparecen estimaciones atípicas cerca de la ra\'iz. Por ello, no afirmamos que la corrida haya verificado empíricamente el orden te\'orico de M1 y M2; harían falta historiales más extensos antes de la saturación numérica y un an\'alisis de un espectro más amplio de semillas.

\section{Extensi\'on a sistemas de ecuaciones no lineales (Semana 9)}
En consonancia con los contenidos de la Semana 9 del curso (Sesión 18) y la bibliograf\'ia de Turner et al. (Cap.~5) \cite{turner}, la resoluci\'on de problemas no lineales en m\'ultiples variables $\mathbf{F}(\mathbf{x}) = \mathbf{0}$, donde $\mathbf{F}: \mathbb{R}^n \to \mathbb{R}^n$, generaliza la derivada escalar a la matriz Jacobiana $\mathbf{J}(\mathbf{x}) \in \mathbb{R}^{n \times n}$, definida por $J_{ij}(\mathbf{x}) = \frac{\partial F_i}{\partial x_j}(\mathbf{x})$. El m\'etodo de Newton multivariable calcula en cada paso el vector de correcci\'on $\mathbf{s}_k$ resolviendo el sistema lineal $\mathbf{J}(\mathbf{x}_k)\mathbf{s}_k = -\mathbf{F}(\mathbf{x}_k)$ mediante descomposici\'on LU con pivoteo parcial ($O(n^3)$ operaciones de punto flotante) para actualizar $\mathbf{x}_{k+1} = \mathbf{x}_k + \mathbf{s}_k$.

Esta extensi\'on dimensional evidencia las limitaciones computacionales del uso de derivadas: ensamblar $\mathbf{J}(\mathbf{x})$ exige evaluar $n^2$ funciones escalares en cada iteraci\'on y expone al algoritmo a divergencias por mal condicionamiento si $\det(\mathbf{J})\approx 0$. Los m\'etodos libres de derivadas de orden superior, como la familia \'optima estudiada en este trabajo, ofrecen una alternativa te\'orica y computacional prometedora al extenderse a sistemas mediante operadores en diferencias divididas multidimensionales o t\'ecnicas de actualizaci\'on tipo Broyden, eliminando la necesidad de derivadas anal\'iticas o c\'alculos jacobianos expl\'icitos.

\section{Alcance, aplicaciones en ciencia de datos y optimizaci\'on}
La Tabla 1 del artículo [1] reporta para g1 con semilla 1.9 siete iteraciones para M1 y M2 y seis para KK, con ACOC 4. En esta corrida M1 converge en cinco pasos, pero M2 y KK se estancan. No es una comparación directa: el artículo usa tolerancia cercana a 1e-200 y 4000 dígitos de precisión en MATLAB; aquí se usan 1e-25 y 150 dígitos. Además, nuestro criterio exige que el paso y el residuo en el nuevo punto cumplan cada uno la tolerancia, mientras que el artículo controla la suma del paso y el residuo del iterado anterior. El menor conteo de M1 no demuestra mayor rapidez si los métodos no cumplen el mismo criterio. Tampoco observamos confirmación del orden cuatro. Las funciones son benchmarks matemáticos deterministas, no datos experimentales.

\textbf{Conexi\'on con Optimizaci\'on y Ciencia de Datos.} El objetivo central de la ciencia de datos en optimizaci\'on radica en minimizar funciones de p\'erdida $L(\boldsymbol{\theta})$ determinando puntos cr\'iticos donde $\nabla L(\boldsymbol{\theta}) = \mathbf{0}$. En paisajes de p\'erdida complejos (como redes neuronales o regresi\'on log\'istica mal condicionada), abundan regiones degeneradas donde la matriz Hessiana $\mathbf{H}(\boldsymbol{\theta})$ es semidefinida positiva con m\'ultiples autovalores nulos; esto corresponde exactamente a puntos estacionarios con multiplicidad $m > 1$. En tales superficies planas, el descenso de gradiente y Newton est\'andar pierden velocidad colapsando a tasas lineales. Los algoritmos de orden \'optimo dise\~nados para multiplicidad conocida permiten acelerar dr\'asticamente la convergencia y superar zonas de estancamiento num\'erico, siendo indispensables para el dise\~no de optimizadores de segundo orden de alto rendimiento.

\section{Cierre e integridad acad\'emica}
En la preparación de este cuaderno y del informe se utilizó una herramienta de IA como apoyo para revisar el plan, organizar el código y mejorar la redacción. La guía del proyecto \cite{guia} permite este apoyo si se declara y no reemplaza la comprensión ni la implementación del equipo. Los integrantes deben verificar las fórmulas, revisar los resultados y poder sustentar el trabajo.
\begin{thebibliography}{9}\small
\bibitem{jerezano} J. H. Jerezano, F. I. Chicharro y N. Garrido-S\'aez, ``Derivative-free optimal multiple-root finding family of iterative methods'', \emph{Computational and Applied Mathematics}, vol. 45, art. 304, 2026. doi: \href{https://doi.org/10.1007/s40314-026-03702-8}{10.1007/s40314-026-03702-8}.
\bibitem{turner} P. R. Turner, T. Arildsen y K. Kavanagh, \emph{Applied Scientific Computing: With Python}, Springer, Cham, 2018. doi: \href{https://doi.org/10.1007/978-3-319-89575-8}{10.1007/978-3-319-89575-8}.
\bibitem{guia} Universidad Externado de Colombia, Departamento de Matemáticas, \emph{Guía del Proyecto Integrador: M\'etodos Numéricos}, 2026.
\end{thebibliography}
\end{document}
